\documentclass[11pt,a4paper]{article}
% Estilo común. Editar aquí para mantener coherentes los tres documentos.
% Adaptación de las plantillas del Seminario Agentes Inteligentes y LLM (UNLu).
\usepackage[utf8]{inputenc}
\usepackage[T1]{fontenc}
\usepackage[spanish,es-nodecimaldot]{babel}
\usepackage[default]{sourcesans}
\usepackage{microtype}
\usepackage[a4paper,left=23mm,right=23mm,top=23mm,bottom=23mm,headheight=15pt,headsep=8mm,footskip=11mm]{geometry}
\usepackage{xcolor}
\usepackage{graphicx}
\usepackage{enumitem}
\usepackage{titlesec}
\usepackage{fancyhdr}
\usepackage{needspace}
\usepackage[skins,breakable]{tcolorbox}
\usepackage{hyperref}
\usepackage{xurl}

\definecolor{azul}{HTML}{1A465E}
\definecolor{petroleo}{HTML}{2F859A}
\definecolor{texto}{HTML}{294D60}
\definecolor{suave}{HTML}{EFF7F9}
\definecolor{borde}{HTML}{CFE2E8}
\definecolor{gris}{HTML}{586D79}

\hypersetup{colorlinks=true,urlcolor=petroleo,linkcolor=azul,citecolor=azul,
  pdfauthor={Aprendizaje Automático - Universidad Nacional de Luján}}
\urlstyle{same}
\color{texto}
\setlength{\parindent}{0pt}
\setlength{\parskip}{6pt}
\linespread{1.10}
\setlength{\emergencystretch}{1.5em}
\clubpenalty=10000
\widowpenalty=10000

% Más aire antes del título que después: el título acompaña su contenido.
\titleformat{\section}{\Large\bfseries\color{azul}}{}{0pt}{}
\titlespacing*{\section}{0pt}{24pt}{0pt}
\titleformat{\subsection}{\large\bfseries\color{azul}}{}{0pt}{}
\titlespacing*{\subsection}{0pt}{18pt}{0pt}
\setlist[itemize]{leftmargin=16pt,itemsep=4pt,topsep=4pt,parsep=0pt,label=\textcolor{petroleo}{\textbullet}}
\setlist[enumerate]{leftmargin=20pt,itemsep=6pt,topsep=4pt,parsep=0pt,label=\textcolor{petroleo}{\bfseries\arabic*.}}

\pagestyle{fancy}
\fancyhf{}
\fancyhead[L]{\small\color{gris}Universidad Nacional de Luján}
\fancyhead[R]{\small\color{gris}Aprendizaje Automático · 2026}
\fancyfoot[L]{\small\color{gris}\pieDocumento}
\fancyfoot[R]{\small\color{petroleo}\thepage}
\renewcommand{\headrulewidth}{0.35pt}
\renewcommand{\footrulewidth}{0pt}
\newcommand{\pieDocumento}{Trabajo práctico integrador}

\newcommand{\cabecera}[3]{%
  {\small\bfseries\color{petroleo}\MakeUppercase{#1}\par}
  \vspace{4pt}
  {\fontsize{24}{28}\selectfont\bfseries\color{azul}#2\par}
  \vspace{7pt}
  {\large\color{gris}#3\par}
  \vspace{6pt}}

\newtcolorbox{destacado}[1]{enhanced,breakable,colback=suave,colframe=borde,
  boxrule=0.5pt,arc=2mm,left=12pt,right=12pt,top=9pt,bottom=9pt,
  borderline west={2pt}{0pt}{petroleo},before skip=12pt,after skip=10pt,
  title={#1},colbacktitle=suave,coltitle=azul,fonttitle=\bfseries,
  attach title to upper={\par\smallskip},fontupper=\normalsize}

\newcommand{\orientacion}[1]{{\color{gris}#1}\par}
\newcommand{\campo}[1]{\textcolor{petroleo}{\textit{[#1]}}}
% Encabezado compacto para los ejemplos consecutivos de la consigna.
\newcommand{\lineadesarrollo}[1]{\par\vspace{12pt}\noindent{\large\bfseries\color{azul}#1\par}\nopagebreak\vspace{-2pt}}

\renewcommand{\pieDocumento}{Plantilla de propuesta · Instancia opcional}
\hypersetup{pdftitle={Plantilla de propuesta - Aprendizaje Automático}}

% Reemplacen los campos y el texto de orientación por su propuesta.
% Esta plantilla adapta la propuesta del Seminario Agentes Inteligentes y LLM.
\begin{document}
\cabecera{Trabajo práctico integrador}{Propuesta de trabajo}{Aprendizaje Automático · Universidad Nacional de Luján}

\textbf{Título tentativo:} \campo{Título del proyecto}\par
\textbf{Estudiante/s:} \campo{Nombre y apellido}\par

\begin{destacado}{Una idea breve para recibir orientación}
Esta instancia es opcional, no tiene fecha de entrega obligatoria y no es requisito para presentar el trabajo final. Escriban una propuesta de \textbf{hasta 200 palabras en total}. Las preguntas siguientes orientan la escritura; pueden responderlas en uno o dos párrafos integrados.
\end{destacado}

\section*{Situación, datos y solución prevista}
\orientacion{¿Qué problema real quieren abordar y para quién sería útil la solución? ¿Qué datos reales utilizarán y de dónde provienen? ¿Qué enfoque de IA piensan aplicar y qué resultado, recomendación o decisión esperan obtener?}

\campo{Escriban aquí la propuesta. Delimiten una tarea realizable y mencionen la fuente de los datos. No hace falta que la solución esté implementada.}

\section*{Relación con la asignatura}
\orientacion{Indiquen brevemente qué conceptos o técnicas del curso sostienen la idea y por qué resultan adecuados. Integren esta explicación dentro del límite de 200 palabras.}

\campo{Conceptos o técnicas de Aprendizaje Automático que se pondrán en práctica.}

\section*{Fuente de los datos}
\campo{Enlace o identificación precisa de la fuente real y forma de acceso.}

\vfill
{\small\color{gris}Reemplacen los campos entre corchetes y retiren las orientaciones antes de enviar. La plantilla organiza la propuesta; también pueden escribirla directamente en la actividad del campus.}
\end{document}
